\chapter{Tensor Algebras, Exterior Algebras, Symmetric Algebras}

Recall the exponential notation introduced in Chapter I, of which we shall make frequent use (I, § 7, no. 8):

Let \((x_{\lambda})_{\lambda \in L}\) be a family of pairwise permutable elements of a ring \(\mathbf{A}\) ; for every mapping \(\alpha: \mathbf{L} \to \mathbf{N}\) of finite support we shall write

\[
x ^ {\alpha} = \prod_ {\lambda \in L} x _ {\lambda} ^ {\alpha (\lambda)}.
\]

If \(\beta\) is another mapping of \(\mathbf{L}\) into \(\mathbf{N}\) of finite support, \(\alpha + \beta\) denotes the mapping

\[
\lambda \mapsto \alpha (\lambda) + \beta (\lambda)
\]

of L into N; with this law of composition the set \(\mathbf{N}^{(L)}\) of mappings of L into N of finite support is the free commutative monoid derived from L and

\[
x ^ {\alpha} x ^ {\beta} = x ^ {\alpha + \beta},
\]

For all \(\alpha \in \mathbf{N}^{(L)}\) , we write \(|\alpha| = \sum_{\lambda \in L} \alpha(\lambda) \in \mathbf{N}\) ; then \(|\alpha + \beta| = |\alpha| + |\beta|\) ; \(|\alpha|\) is called the order of the ``multiindex'' \(\alpha\) . For all \(\lambda \in L\) , let \(\delta_{\lambda}\) denote the element of \(\mathbf{N}^{(L)}\) such that \(\delta_{\lambda}(\lambda) = 1\) , \(\delta_{\lambda}(\mu) = 0\) for \(\mu \neq \lambda\) (Kronecker index); the \(\delta_{\lambda}\) for \(\lambda \in L\) are the only elements of \(\mathbf{N}^{(L)}\) of order 1. \(\mathbf{N}^{(L)}\) is given the ordering induced by the product ordering on \(\mathbf{N}^{\mathrm{L}}\) , so that the relation \(\alpha \leqslant \beta\) is equivalent to `` \(\alpha(\lambda) \leqslant \beta(\lambda)\) for all \(\lambda \in L\) ''; then the multiindex \(\lambda \mapsto \beta(\lambda) - \alpha(\lambda)\) is denoted by \(\beta - \alpha\) , so that it is the unique multiindex such that \(\alpha + (\beta - \alpha) = \beta\) . For all \(\alpha \in \mathbf{N}^{(L)}\) , there are only a finite number of multiindices \(\beta \leqslant \alpha\) ; the \(\delta_{\lambda}\) are the minimal elements of the set \(\mathbf{N}^{(L)} = \{0\}\) ; the relation \(\alpha \leqslant \beta\) implies \(|\alpha| \leqslant |\beta|\) and if both \(\alpha \leqslant \beta\) and \(|\alpha| = |\beta|\) , then \(\alpha = \beta\) . Finally, we write \(\alpha! = \prod_{\lambda \in L} (\alpha(\lambda))!\) , which is meaningful since \(0! = 1\) .

From § 4 to § 8 inclusive, A denotes a commutative ring and, unless otherwise stated, the algebras considered are assumed to be associative and unital and the algebra homomorphisms are assumed to be unital.

